\label{prob:medical-test}
A medical test finds a serious disease in a patient. The test is 99\%
accurate: if the patient has the disease, the test is positive with
probability 99\%; if the patient does not have the disease, the test is
negative with probability 99\%. The disease is rare: only one person in
10\,000 has it. Find the probability that the patient really has the
disease.

\begin{solution}
\solpart{Interpretation}
Let $D \in \{0, 1\}$ show whether the patient has the disease ($D = 1$).
Let $T \in \{0, 1\}$ be the test result ($T = 1$ is a positive result).
The problem gives three values:
\[
  p(D = 1) = 10^{-4}, \qquad
  p(T = 1 \mid D = 1) = 0.99, \qquad
  p(T = 0 \mid D = 0) = 0.99.
\]
Find $p(D = 1 \mid T = 1)$.

\solpart{Idea}
The disease causes the test result, so the causal graph is $D \to T$.
The graph gives the joint distribution as a product of two factors:
\[
  p(D, T) = p(D) \, p(T \mid D).
\]
The problem gives both factors: the distribution of the source node $D$
and the conditional distribution on the edge $D \to T$. Thus the joint
distribution is known. From the joint distribution, you can get any other
probability with two operations:
\begin{itemize}
  \item A marginal probability is the sum of the joint probability over
    the other variables. For a continuous variable, the sum is an
    integral.
  \item A conditional probability is a joint probability divided by a
    marginal probability.
\end{itemize}
These two operations together give Bayes' theorem.

\solpart{Solution}
\solstep{Step 1. Find the joint probabilities of a positive result.}
The probability of a false positive result is
$p(T = 1 \mid D = 0) = 1 - 0.99 = 0.01$. Thus
\begin{align*}
  p(D = 1, T = 1) &= p(D = 1) \, p(T = 1 \mid D = 1)
    = 10^{-4} \cdot 0.99 = 0.000099, \\
  p(D = 0, T = 1) &= p(D = 0) \, p(T = 1 \mid D = 0)
    = 0.9999 \cdot 0.01 = 0.009999.
\end{align*}

\solstep{Step 2. Find the marginal probability of a positive result.}
Sum the joint probabilities over $D$:
\[
  p(T = 1) = 0.000099 + 0.009999 = 0.010098.
\]

\solstep{Step 3. Divide the joint probability by the marginal probability.}
\[
  p(D = 1 \mid T = 1) = \frac{p(D = 1, T = 1)}{p(T = 1)}
  = \frac{0.000099}{0.010098} = \frac{1}{102} \approx 0.0098.
\]
The patient has the disease with a probability of approximately 1\%.

\solpart{Conclusion}
A positive result increases the probability of the disease from 0.01\% to
approximately 1\%. But the patient is still healthy with a probability of
99\%. The cause is the low prior probability. In 10\,000 people, the test
finds approximately 1 sick person and 100 healthy people. For this reason,
most positive results of a mass test for a rare disease are false. A second
independent test gives more information: if it is also positive, the
probability of the disease increases to 0.495.
\end{solution}
